\subsection{Interpretability of the Learned LGCP Parameters}
\label{app:interpretability}

Unlike the neural baselines, the LGCP has a small set of learned parameters with a direct statistical meaning. The log-intensity of a cell $\mathbf{s}$ on day $t$ is
\begin{equation}
\log\lambda(\mathbf{s},t)=\alpha+\boldsymbol{\beta}^\top\mathbf{z}(\mathbf{s},t)+f(\mathbf{s},t),
\qquad
f\sim\mathcal{GP}\bigl(0,\,(k_{s,0}+k_{s,1})(k_{t,0}+k_{t,1})\bigr),
\label{eq:lgcp_interp}
\end{equation}
where $\mathbf{z}$ collects the eight covariates, $k_{s,0}$ and $k_{s,1}$ are RBF kernels on the standardized cell coordinates with lengthscales $0.2$ and $0.01$, $k_{t,0}$ is an RBF kernel on standardized time with lengthscale $0.01$, and $k_{t,1}$ is a periodic kernel with lengthscale $0.8$. The lengthscales are fixed, while the intercept $\alpha$, the coefficients $\boldsymbol{\beta}$, the four kernel variances $\sigma^2_{s,0},\sigma^2_{s,1},\sigma^2_{t,0},\sigma^2_{t,1}$ and the period $p$ are optimized. Table~\ref{tab:lgcp_params} reports their values as mean $\pm$ standard deviation over the 10 seeds $\times$ 4 weekly windows of each month. The zero-inflation gate is applied on top of $\lambda$ and is not described by these parameters.

\begin{table}[t]
\centering
\caption{Optimized LGCP parameters (mean $\pm$ population standard deviation over 10 seeds $\times$ 4 weekly windows per month). The ``Effect'' columns give the multiplicative effect on the expected count, $\exp(\beta/\mathrm{sd})$, per physical unit or for switching a binary indicator on, computed per run with its own training statistics (for chlorophyll-a, $\exp(\beta)$ per training standard deviation); for $\alpha$ they give the effective intercept $\alpha+\mathrm{mean}(\bar f)$, and for the kernel rows the fixed lengthscales in physical units. The bottom block reports scale-free kernel summaries and the share of variance of the posterior-mean field $\bar f$, over all training cell-days, carried by each kernel component; the shares need not sum to 100\% because the components are correlated.}
\label{tab:lgcp_params}
\vspace{2pt}
\footnotesize
\setlength{\tabcolsep}{3.3pt}
\begin{tabular}{llccccl}
\toprule
 & & \multicolumn{2}{c}{Optimized value} & \multicolumn{2}{c}{Effect} & \\
\cmidrule(lr){3-4}\cmidrule(lr){5-6}
Parameter & Description & May & Nov. & May & Nov. & Effect unit \\
\midrule
\multicolumn{7}{l}{\textit{Linear predictor} $\log\lambda=\alpha+\boldsymbol{\beta}^\top\mathbf{z}+f$} \\
$\alpha$ & Intercept & $-2.129 \pm 0.013$ & $-2.098 \pm 0.013$ & $-2.982$ & $-2.887$ & $\alpha+\mathrm{mean}(\bar f)$ \\
$\beta_{\mathrm{chl}}$ & Chlorophyll-a & $0.079 \pm 0.007$ & $0.092 \pm 0.018$ & $\times1.082$ & $\times1.097$ & per $+1$ SD \\
$\beta_{\mathrm{temp}}$ & Sea temperature & $0.324 \pm 0.009$ & $0.342 \pm 0.013$ & $\times1.052$ & $\times1.054$ & per $+1\,^\circ$C \\
$\beta_{\mathrm{sal}}$ & Salinity & $0.003 \pm 0.011$ & $0.042 \pm 0.009$ & $\times1.003$ & $\times1.034$ & per $+1$ psu \\
$\beta_{\mathrm{ban}}$ & Fishing ban & $-0.918 \pm 0.011$ & $-1.027 \pm 0.011$ & $\times0.041$ & $\times0.048$ & on vs.\ off \\
$\beta_{\mathrm{hol}}$ & Public holiday & $-0.400 \pm 0.009$ & $-0.341 \pm 0.011$ & $\times0.122$ & $\times0.148$ & on vs.\ off \\
$\beta_{\mathrm{wkend}}$ & Weekend & $-0.719 \pm 0.012$ & $-0.662 \pm 0.012$ & $\times0.203$ & $\times0.231$ & on vs.\ off \\
$\beta_{\mathrm{lag}1}$ & Count at $t-1$ & $0.191 \pm 0.004$ & $0.203 \pm 0.006$ & $\times1.234$ & $\times1.234$ & per $+1$ count \\
$\beta_{\mathrm{lag}7}$ & Count at $t-7$ & $0.068 \pm 0.006$ & $0.075 \pm 0.004$ & $\times1.078$ & $\times1.082$ & per $+1$ count \\
\midrule
\multicolumn{7}{l}{\textit{Kernel} $k=(k_{s,0}+k_{s,1})(k_{t,0}+k_{t,1})$, fixed lengthscales $\ell$} \\
$\sigma^2_{s,0}$ & Spatial RBF, $\ell{=}0.2$ & $1.987 \pm 0.179$ & $1.608 \pm 0.122$ & 2.8\,/\,1.4\,km & 2.8\,/\,1.4\,km & $\ell$ E--W\,/\,N--S \\
$\sigma^2_{s,1}$ & Spatial RBF, $\ell{=}0.01$ & $0.088 \pm 0.008$ & $0.101 \pm 0.008$ & 0.14\,/\,0.07\,km & 0.14\,/\,0.07\,km & $\ell$ E--W\,/\,N--S \\
$\sigma^2_{t,0}$ & Temporal RBF, $\ell{=}0.01$ & $0.100 \pm 0.006$ & $0.096 \pm 0.005$ & 1.4\,d & 2.0\,d & $\ell$ in days \\
$\sigma^2_{t,1}$ & Periodic, $\ell{=}0.8$ & $0.873 \pm 0.054$ & $0.880 \pm 0.046$ &  &  &  \\
$p$ & Period (days), init.\ 14 & $14.03 \pm 0.01$ & $14.02 \pm 0.07$ &  &  &  \\
\midrule
\multicolumn{7}{l}{\textit{Scale-free kernel summaries and variance decomposition of} $\bar f$} \\
Prior var. & $(\sigma^2_{s,0}{+}\sigma^2_{s,1})(\sigma^2_{t,0}{+}\sigma^2_{t,1})$ & $2.01 \pm 0.09$ & $1.66 \pm 0.07$ &  &  &  \\
Fine share & $\sigma^2_{s,1}/(\sigma^2_{s,0}{+}\sigma^2_{s,1})$ & 4.3\% & 5.9\% &  &  &  \\
Periodic share & $\sigma^2_{t,1}/(\sigma^2_{t,0}{+}\sigma^2_{t,1})$ & 89.7\% & 90.2\% &  &  &  \\
Static & Time-constant cell map & 93.7\% & 94.3\% &  &  &  \\
Cycle & Periodic (weekly/fortnightly) part & 5.6\% & 5.2\% &  &  &  \\
Short-term & Temporal RBF part & ${<}0.1\%$ & ${<}0.1\%$ &  &  &  \\
\bottomrule
\end{tabular}
\end{table}

\paragraph{From parameters to effects.}
All covariates, binary indicators included, are standardized with the mean and standard deviation of the training set. A coefficient $\beta_j$ in~\eqref{eq:lgcp_interp} is therefore the change in log-intensity per training standard deviation $\mathrm{sd}_j$ of covariate $j$, and the multiplicative effect on the expected count of a one-unit increase of the raw covariate, or of switching a binary indicator from 0 to 1, is $\exp(\beta_j/\mathrm{sd}_j)$; for an indicator active on a fraction $q_j$ of the training days, $\mathrm{sd}_j=\sqrt{q_j(1-q_j)}$. We apply this conversion separately to every run, using its own training statistics, and then average. This matters because the training statistics differ between windows: for instance, the 2025 fishing ban enters the training data of the November windows only, which raises the standard deviation of the ban indicator from $0.29$ to $0.34$, so raw coefficients should not be compared across months. The intercept on its own is not a baseline rate: $e^{\alpha}$ is the intensity at average covariate values and $f=0$, but the posterior mean $\bar f$ of the latent field is not centered (its average over the training cell-days is $-0.85$ in May and $-0.79$ in November), so the effective intercept is $\alpha+\mathrm{mean}(\bar f)\approx-2.98$ and $-2.89$, respectively.

\paragraph{Identifiable kernel quantities.}
Because the spatial and temporal kernels are multiplied, rescaling both spatial variances by any $c>0$ and both temporal variances by $1/c$ leaves the covariance, and hence the ELBO, unchanged. The individual variances are therefore determined only up to this factor, and we interpret the kernel through scale-free quantities: the within-group shares $\sigma^2_{s,1}/(\sigma^2_{s,0}+\sigma^2_{s,1})$ and $\sigma^2_{t,1}/(\sigma^2_{t,0}+\sigma^2_{t,1})$, and the product $(\sigma^2_{s,0}+\sigma^2_{s,1})(\sigma^2_{t,0}+\sigma^2_{t,1})$, which is the prior variance of $f$ at any point. In addition, we decompose the posterior mean $\bar f(\cdot)=K_{\cdot Z}K_{ZZ}^{-1}\mathbf{m}$, where $Z$ are the inducing inputs and $\mathbf{m}$ the variational mean. Expanding the product kernel into its four cross terms, and splitting the periodic kernel into its time-constant Fourier term and an oscillating remainder,\footnote{With lengthscale $0.8$, the periodic kernel $k_{t,1}(\tau)\propto\exp\bigl(-2\sin^2(\pi\tau/p)/\ell^2\bigr)$ places 36\% of its variance on a constant term, 44\% on period $p$, 16\% on period $p/2$ and about 5\% on higher harmonics.} writes $\bar f$ as the sum of a time-constant spatial map, a periodic component and a short-term component. The variance of each component over all training cell-days, relative to the variance of $\bar f$, gives the shares in the bottom block of Table~\ref{tab:lgcp_params}.

\paragraph{Covariate effects.}
Calendar and regulatory factors have by far the largest effects. During the fishing ban (31~July--13~September) the expected count falls to about 4--5\% of its value outside the ban, on public holidays to 12--15\%, and on weekends to 20--23\%. These effects are consistent with the raw data, where the corresponding ratios of mean daily counts are 0.015--0.019 (ban vs.\ non-ban weekdays in June--October), 0.06 (holidays vs.\ other weekdays, pooled over 2024--2025) and 0.16 (weekend days vs.\ weekdays); the model's effects are smaller in magnitude because they are conditional on the lagged counts and on $f$, which absorb part of these drops. Activity is also persistent: each additional count observed in the same cell on the previous day multiplies the expected count by $1.23$, and each count observed seven days earlier by $1.08$. Among the environmental covariates, sea temperature has a positive effect of about 5\%~per~$^\circ$C ($\times 1.38$--$1.41$ per training standard deviation of $6.4\,^\circ$C); as temperature varies mostly in time rather than across the gulf, this effect largely reflects higher activity in the warm season, net of the summer fishing ban. Chlorophyll-a has a positive effect in all 80 runs ($\times 1.08$--$1.10$ per standard deviation): moving from the 5th to the 95th percentile of its training distribution (about $0.18$ to $1.04$\,mg\,m$^{-3}$) raises the expected count by 29\% in May and 34\% in November, consistent with vessel activity tracking primary productivity. Because chlorophyll-a varies much more in time than across the gulf, this effect is driven mainly by its temporal variation. The daily chlorophyll-a product also exhibits a weekly sawtooth pattern, with a jump at the start of each week followed by a gradual decay (the day of the week explains about 80\% of the variance of daily chlorophyll-a around its weekly running mean; the pattern is already present in the source Copernicus products), so we verified that the effect is not a proxy for the day of the week or for seasonal and inter-annual trends. In a Poisson GLM with cell-specific intercepts and the same covariates, fitted on the training data of the first May and November windows, the chlorophyll-a coefficient is positive in all 16 specifications we considered, with or without day-of-week, year and month indicators and using either daily or 7-day-averaged chlorophyll-a ($0.05$--$0.19$ per standard deviation), and it is significant at the 5\% level in 15 of them (quasi-Poisson likelihood-ratio test). Without additional indicators, the same GLM also recovers the calendar, temperature and lag coefficients of the LGCP closely (e.g., in May, fishing ban $-1.08$ vs.\ $-0.92$, weekend $-0.72$ vs.\ $-0.72$, temperature $0.38$ vs.\ $0.32$, lag-1 count $0.17$ vs.\ $0.19$), indicating that these effects are not an artifact of the GP. Salinity has a smaller effect, which is positive only in the November runs: there, its coefficient is positive in all 40 runs, and moving from the 5th to the 95th percentile of its training distribution (about $34$ to $38$\,psu) raises the expected count by 15\%, whereas in May it is close to zero (positive in 22 of 40 runs). This difference is consistent with the salinity coefficient partly reflecting the contrast between the two years: salinity was on average about 1\,psu higher in 2025 than in 2024, 2025 was also the busier year in the November training data, and in the GLM the salinity coefficient is not significant once a year indicator is included ($p\ge0.31$).

\paragraph{Spatial structure.}
Because the coordinates are standardized separately along each axis, the lengthscale $0.2$ of $k_{s,0}$ corresponds to about $2.8$\,km in the east--west and $1.4$\,km in the north--south direction, which is below the cell size ($3.25 \times 4.63$\,km). Under $k_{s,0}$, neighboring cells are correlated at $0.51$ in the east--west direction and are essentially uncorrelated in the north--south direction, while $k_{s,1}$ (lengthscale $0.01$) is independent across cells. The spatial part of the GP therefore acts as a cell-level random effect with mild east--west smoothing, rather than as a smooth large-scale surface. The decomposition confirms this: about 94\% of the variance of $\bar f$ is a time-constant map of cell effects. This map closely tracks the mean activity of each cell (Spearman correlation $0.97$--$0.98$ with the mean training count), is virtually identical in May and November (correlation $0.99$), and spans a factor of about 25--30 in expected intensity between the most and the least active cells. The fine component $k_{s,1}$ shrinks from its initial variance of $0.5$ to about $0.09$--$0.10$ and accounts for only 4--6\% of the spatial prior variance.

\paragraph{Temporal structure.}
The periodic kernel carries about 90\% of the temporal prior variance: during training its variance grows from $0.2$ to $0.87$--$0.88$, while that of the short-range RBF kernel $k_{t,0}$ decreases from $0.5$ to about $0.10$. The learned period remains at approximately 14 days in all runs ($13.76$--$14.28$). With the fixed lengthscale of $0.8$, a 14-day periodic kernel also spans the 7-day harmonic, and the learned periodic component of $\bar f$ is in fact predominantly weekly: after averaging over cells and over the runs of each month, 88--94\% of its variance over the 14-day cycle repeats with a 7-day period. The resulting weekly profile is highest on Sundays and Mondays and lowest on Wednesdays and Thursdays, with a peak-to-trough ratio of about $1.2$ in expected intensity. Since the weekend indicator and the lagged counts already model the main weekly drop, this component acts as a residual day-of-week correction, capturing for instance that Sundays are busier than Saturdays and that activity recovers on Mondays, which a single weekend indicator and the lag-1 term cannot represent. Short-term fluctuations, finally, are captured by the lagged counts rather than by the GP: the short-range RBF kernel, whose lengthscale corresponds to $1.4$--$2.0$ days depending on the window, accounts on average for less than 0.1\% of the variance of $\bar f$, as variation on the scale of a single cell and a few days cannot be resolved with 300 inducing points.

\paragraph{Stability.}
The calendar, temperature and lag coefficients are stable across seeds and windows: within each month, their relative standard deviation is below 4\% (below 9\% for the lag-7 coefficient), and the sign of every coefficient except salinity is the same in all 80 runs. The chlorophyll-a and salinity coefficients vary more across runs (Table~\ref{tab:lgcp_params}). These effects describe conditional associations learned from observational data rather than causal effects; in particular, the three calendar indicators are constant across cells on a given day and can partly trade off with the temporal component of $f$.
